    % Escreva aqui a solução da questão 2.
\textbf{VERDADEIRA}

\textcolor{red}{ VERSÃO 1 } %Letícia

De uma das propriedades da interseção de conjuntos, temos que $A \cap B^{C} \subseteq B^{C}$.
Então, podemos reescrever:
\begin{align*}
    A \cap B^{C} \subseteq B^{C} &\implies (B^{C})^{C} \subseteq (A \cap B^{C})^{C} \tag{$X \subseteq Y \implies Y^C \subseteq X^C$} \\ 
    & \implies B \subseteq (A \cap B^{C})^{C} \tag{$(X^C)^C = X$}\\ 
    &\implies B \subseteq (A \setminus B)^{C} . \tag{$X\setminus Y  = X \cap Y^{C}$}
\end{align*}

Logo, a relação de inclusão é verdadeira.




\bigskip

\textcolor{red}{CHAVE DE PONTUAÇÃO:
\begin{itemize}
    \item  Inclusão pela interseção \emph{(1,0 ponto)};
    \item Uso de $X \subseteq Y \implies Y^C \subseteq X^C$ e do complementar do complementar\emph{(1,0 ponto)};
    \item  Uso da igualdade $X \setminus Y = X \cap Y^C$ \emph{(1,0 ponto)}.
\end{itemize} } 

\bigskip


\textcolor{red}{ VERSÃO 2} %Letícia


Analisando o conjunto do lado direito da relação de inclusão, perceba que a diferença entre os conjuntos pode ser reescrita como a interseção do conjunto à esquerda da diferença ($A$) e do complementar do conjunto à direita ($B$), logo 
\begin{align}
    (A\setminus B)^{C} & = (A \cap B^{C})^{C} & \tag{diferença de conjuntos} \\
    & = A^{C} \cup B. & \tag{De Morgan e $(X^C)^C = X$}
\end{align}

Então,  $B \subseteq (A \setminus B)^{C}$ é equivalente a $B \subseteq A^{C} \cup B$. De uma das propriedades da união de conjuntos, sabe-se que essa relação de inclusão é verdadeira.



\bigskip

%\textcolor{red}{ VERSÃO 3} % autor




%\bigskip

\textcolor{red}{CHAVE DE PONTUAÇÃO:
\begin{itemize}
    \item  Inclusão pela união \emph{(1,0 ponto)};
    \item  Igualdade (ou inclusão) pelo DeMorgan e o complementar do complementar\emph{(1,0 ponto)};
    \item  Uso da igualdade $X \setminus Y = X \cap Y^C$ \emph{(1,0 ponto)}.
\end{itemize} }